\section{Код в Wolfram Mathematica для решения уравнения осцилляции нейтрино в среде}
\label{app:code-python-vis}

Для численого решения уравнения \eqref{eq:15} использовалась система компьютерной алгебры \textsf{Wolfram Mathematica}. 


Ниже представлен код сценария, использованного для получения данных по квазипериодам \(Re \Psi_1(\xi)\), \(Im\Psi_1(\xi)\) в зависимости от \(\xi\), а также количества переходов через ось обцис \(\Re\Psi_{2,3}\) и \(\Im\Psi_{2,3}\) на интервале каждого из ранее наайденых квазипериодов \(\Re\Psi_1\) и \(\Im\Psi_1\).

Его также можно найти в репозитории \url{https://github.com/Fest-Fut/Kursovay}.

\inputminted[linenos,breaklines,fontsize=\scriptsize]{Mathematica}{../Mathematica/Sun_DOPRI_p5.nb}

%\textattachfile{stat-prop_v2.py}{Код в виде прикреплённых данных (только для
%	электронной версии документа).}
